\section{Related Work}
\label{sec:related}

A chain whose block validity depends on external layers raises five
questions, and each has been studied by a separate community.
Fig.~\ref{fig:taxonomy} organizes the literature by these questions:
how consensus behaves under degraded conditions (Section~\ref{subsec:rw-degraded}),
how local observations of an external layer become data that all validators
see and agree on (Section~\ref{subsec:rw-detect}), how modular designs cope with
the failure of Bitcoin, DA, or settlement layers (Section~\ref{subsec:settlement}),
how a chain returns to the anchored state (Section~\ref{subsec:rw-recovery}),
and how such designs are validated (Section~\ref{subsec:rw-validation}).
Several building blocks of Engram exist in prior work: vote extensions that carry
signed external data, weighted-median aggregation, and protocol-defined
fallback for DA and sequencing failures. Engram differs in what the aggregated
observation controls. Table~\ref{tab:gates} shows what the failure of a dependency stops in each
design, and Appendix~\ref{app:related} compares representative systems by trigger,
decision authority, behavior under failure, recovery, assumption and reported
assurance; Section~\ref{subsec:rw-gap} states the combination left open.

The taxonomy of related work (Fig.~\ref{fig:taxonomy}) is in Appendix~\ref{app:extended}.

\subsection{Consensus Under Degraded Conditions}
\label{subsec:rw-degraded}
Related work on agreement under degraded conditions (agreement core, fast paths, availability--finality trade-offs) is in Appendix~\ref{app:extended}.

\subsection{From Local Observations to Consensus-Visible External Data}
\label{subsec:rw-detect}
Related work on failure detectors, the aggregation of Byzantine observations, oracles, vote extensions and peripheral P2P attacks is in Appendix~\ref{app:extended}.

\subsection{Fallback for Modular-Chain Dependencies}
\label{subsec:settlement}
\textbf{Data availability.}
LazyLedger/Celestia~\cite{albassam2019lazyledger} decouples DA from execution.
Fraud and data availability proofs~\cite{albassam2021fraud} let light clients
detect withheld data, and coded distributed arrays target the peer-to-peer layer that
distributes data availability sampling fragments~\cite{phamminh2026cda}. This line of work makes DA
\emph{verifiable}, but a chain that requires a DA inclusion proof before
accepting a block still stops when DA stops. Engram adds a consensus-level
response to DA unavailability and a backfill requirement before the chain
returns to \textsc{Anchored} (Section~\ref{subsec:reanchoring}).

\textbf{Bitcoin anchoring.}
Pegged sidechains~\cite{back2014sidechains} introduced two-way pegs to Bitcoin.
Karakostas and Kiayias~\cite{karakostas2021checkpointing} study checkpointing
as a defense of proof-of-work ledgers against 51\% attacks and report a
liveness attack, ``block lead'', against existing checkpointing designs. For
proof-of-stake chains, Babylon~\cite{tas2022babylon,tas2023bitcoinpos} and
Pikachu~\cite{azouvi2022pikachu} checkpoint onto Bitcoin to prevent long-range
attacks. Tas et al.~\cite{tas2023bitcoinpos} show that slashable safety is
impossible without an external source of trust and state their guarantees under
the assumption that Bitcoin is safe and live. Their design also defines a
fallback in which Bitcoin orders the PoS blocks after censorship, which
strengthens the PoS chain's dependence on Bitcoin. Babylon specifies a fallback for censored checkpoints and an emergency break but, to our knowledge, none of
these publications specifies how the chain operates while the anchor itself is
unreachable. Engram treats unreachability as an input backed by a quorum of
signed observations, keeps ordering blocks, and locks withdrawals while it
persists. Accountability comes from CometBFT's equivocation evidence; Engram
implements no Bitcoin-enforced slashing. A rollup that settles its state root on Bitcoin gives users an exit independent of the operators; Engram instead keeps its own validator set sovereign and trades exit for liveness, and we do not claim that the trade is preferable for asset custody.

The comparison with rollup fallback and its decision locus is in Appendix~\ref{app:extended}.

\textbf{Withdrawal risk.}
Lee et al.~\cite{lee2023bridges} break a bridge into its components, identify
the risk each introduces, and analyze past exploits that targeted bridges. They
report that over USD~1 billion in digital assets was stolen, incorrectly
minted, or locked in bridges in the year before their study. Withdrawals are the
path through which a chain without external corroboration can lose assets,
which is why Engram locks them while the chain runs in \textsc{Sovereign} mode.

\subsection{Recovery and Re-Anchoring}
\label{subsec:rw-recovery}
Succinct proofs are the standard tool for convincing a verifier of a chain's
history without re-executing it. They differ in how a long history is covered.
Mina~\cite{bonneau2020mina} keeps a constant-size proof of the whole chain by
recursively composing SNARKs, which Plumo's authors report imposes a large
overhead on the protocol being proved~\cite{vesely2022plumo}. zkBridge~\cite{xie2022zkbridge}
verifies source-chain consensus inside a two-layer recursive proof system.
Plumo instead syncs a client through a chain of sequential SNARK transition
proofs, each covering four months of blockchain history~\cite{vesely2022plumo}.
Recursive composition~\cite{bitansky2013recursive,bowe2019halo} and
folding~\cite{kothapalli2022nova} are the general techniques behind the
recursive designs. FlyClient~\cite{bunz2020flyclient} takes a different route
for proof-of-work chains: a light client downloads only a logarithmic number of
block headers, selected by probabilistic block sampling.
All of these prove that a header chain, or a validator-set transition, is
valid. A chain that has run without external corroboration must also show that
the interval followed the protocol. The long-range attack
literature~\cite{deirmentzoglou2019longrange} reviews how an attacker can
rewrite the history of a proof-of-stake chain, which is the exposure of an
interval secured only by the chain's own validators. Engram's recovery proof is
a fixed-capacity PLONK-family~\cite{gabizon2019plonk} proof over one interval
chunk, and chunk proofs are chained through an on-chain root, as Plumo chains
its transition proofs. The statement differs: besides header continuity it covers
the recorded mode and lock-flag label of every header and the quorum signatures
of the two boundary blocks, and it does not check transitions against the
transition function (Section~\ref{subsec:reanchoring}).

Related work on validation methodology is in Appendix~\ref{app:extended}.

\subsection{Research Gap}
\label{subsec:rw-gap}
Prior work supplies consensus that adapts to internal participation faults,
oracle systems that aggregate external observations, vote extensions that
deliver signed external data to the validator pipeline, rollup fallbacks for
sequencing and DA failures, and succinct proofs of chain histories. When an
external layer fails, a chain that gates validity on it halts, and a chain
that drops the dependency gives up the guarantees that motivated it. The
fallbacks surveyed above place the decision with users, watchers, a sequencer,
a committee, or an L1 contract.

\textbf{What a response gates.} The designs differ in what the failed
dependency stops (Table~\ref{tab:gates}). Only the strict receipt-gated baselines
stop block commitment. The other designs stop finality, client acceptance, L2
derivation or nothing, or leave the chain running and gate an operation.

\begin{table}[!tbp]
\caption{What the failure of an external dependency stops, by design (entries
from the sources of Table~\ref{tab:comparison}, except the baseline and Engram).}
\label{tab:gates}
\centering
\footnotesize
\renewcommand{\arraystretch}{1.03}
\begin{tabular}{@{}p{0.2\linewidth}p{0.4\linewidth}p{0.3\linewidth}@{}}
\toprule
\textbf{Gate} & \textbf{What stops} & \textbf{Design} \\
\midrule
Block validity & Every proposal, while the validity predicate cannot be satisfied & Strict receipt-gated baselines, in-band and out-of-band (this work) \\
Finality & The finalized prefix lags; the ledger keeps growing & Ebb-and-flow \\
Client acceptance & A client stops adding blocks whose data behind a signed checkpoint is unavailable & Babylon emergency break \\
L2 derivation & Derivation stalls until a challenge resolves the data or the blocks are reorged & OP Alt-DA \\
Oracle input & Nothing: an empty vote extension is valid and the chain keeps running & Vote-extension oracle \\
An operation & Withdrawals, and the return to the anchored mode; block production continues & Engram \\
\bottomrule
\end{tabular}
\end{table}

To the best of our knowledge, no published design combines (i)~a quorum-backed
set of height-bound validator observations from which every validator
deterministically derives the health of external layers, (ii)~use of that
health inside the validator consensus pipeline to change both proposal-validity
requirements and withdrawal admissibility while the native validator quorum
keeps ordering blocks, and (iii)~a recovery condition that requires a proof over the degraded
interval and the backfill of that interval to the DA layer before the chain
returns to the externally anchored mode. Engram provides these three (with the proof-gate limitation of
Section~\ref{subsec:limitations}), retains CometBFT as the agreement protocol,
and Sections~\ref{sec:model}--\ref{sec:evaluation} evaluate the combination.


